% Who Wins When One Agent Reflects?
% Pass-Through Effects of Prompt Interventions in Multi-Agent Games
% Augustin Chan, 2026

\documentclass[11pt]{article}
\usepackage[margin=1in]{geometry}
\usepackage{amsmath,amssymb}
\usepackage{graphicx}
\usepackage{booktabs}
\usepackage{hyperref}
\usepackage{natbib}
\usepackage{xcolor}

\title{Who Wins When One Agent Reflects?\\
\large Pass-Through Effects of Prompt Interventions in Multi-Agent Games}
\author{Augustin Chan\\
\texttt{aug@iterative.day}}
\date{May 2026}

\begin{document}
\maketitle

\begin{abstract}
We inject symbolic reasoning frameworks into one agent in a 7-player simultaneous-move strategy game and measure ecosystem-level outcomes across 41 games (4 conditions, $n \geq 10$ per cell). The intervention does not benefit its recipient---survival drops from 64\% (control) to 30--40\% (all framework conditions)---but systematically redirects which \emph{other} state wins. Under yarrow-stalk I~Ching divination, the strongest expansionist state (Qin) wins 0 of 10 games (Fisher $p = 0.033$ vs.\ tarot). Under Tarot, Qin wins 5 of 10. Under scrambled (incoherent) text, the intelligence-school state (Qi) wins 6 of 10 (Fisher $p = 0.040$ vs.\ others). A perturbativeness gradient in non-recipient reasoning length (control~98 $<$ yarrow~126 $<$ tarot~152 $<$ scrambled~197 chars/order; Kruskal-Wallis $p < 0.001$) provides a candidate mechanism: more perturbative frameworks induce longer reasoning in \emph{other} agents. We name this pattern class \textbf{pass-through effects}---interventions that do not improve their target but redirect system-level outcomes to third parties.
\end{abstract}

% ==============================================================
\section{Introduction}
\label{sec:intro}

% TODO: Lead with the punchline. Cite paper 1. Preview structure.

% ==============================================================
\section{Experimental Setup}
\label{sec:setup}

\subsection{Environment and Agents}
% 7-state Warring States diplomacy game, Diplomacy-style mechanics,
% 20-round games, LLM agents (Claude Opus 4.6)

\subsection{Intervention Design}
% Dual-timescale: decision-time (per-round MANDATE) + learning-time (between-game reflection)
% Length-matched control

\subsection{Conditions and Controls}
% 4 conditions: control, yarrow, tarot, scrambled
% Table with description of each

\subsection{Metrics}
% Local: Han survival, peak SCs
% System: winner distribution, elimination order

% ==============================================================
\section{Results I: No Local Benefit}
\label{sec:local}

% Han survival, peak SCs, tarot elevates peak (KW p=0.008)
% But survival worse under all frameworks

% ==============================================================
\section{Results II: Ecosystem Redirection}
\label{sec:ecosystem}

% Winner distributions, Qin suppression, Qi dominance
% "Redistribution, not degradation"

% ==============================================================
\section{Results III: Perturbativeness Gradient}
\label{sec:gradient}

% Non-Han reasoning length monotonic increase
% Dose-response: more perturbative framework → more reasoning disruption

% ==============================================================
\section{Synthesis: Pass-Through Effects}
\label{sec:synthesis}

% Define the pattern class
% Connect sections 3-5
% Diagram: Intervention → Agent A → (no improvement) → System redistribution → Agents B, C...

% ==============================================================
\section{Discussion}
\label{sec:discussion}

% Alignment implications: optimizing agents ≠ optimizing systems
% Evaluation implications: local metrics insufficient
% Generalizability: where else might pass-through effects appear?

% ==============================================================
\section{Limitations}
\label{sec:limitations}

% Single game environment
% Specific prompt class (symbolic/reflective)
% n=10 per cell (power analysis in supplementary)
% Cannot fully ablate decision-time vs memory effects
% Single model (Claude Opus 4.6)

% ==============================================================
\section{Conclusion}
\label{sec:conclusion}

% Modifying one agent need not improve that agent.
% It can instead redirect the system.

\bibliographystyle{plainnat}
\bibliography{references}

\end{document}
